\begin{table}[!htbp]
\centering\small
\setlength{\tabcolsep}{5pt}
\renewcommand{\arraystretch}{1.10}
\caption{TV--persistence sensitivity relative to its ridge reference.}
\label{tab:legacy_sensitivity}
\begin{tabular}{lrr}
\toprule
Variant & Lower MSE & Median PRE (\%) \\
\midrule
Default & 9/15 & +6.72 \\
Sample standard deviation & 9/15 & +5.18 \\
Rousseeuw--Croux $Q_n$ & 9/15 & +6.72 \\
Common scale across horizons & 9/15 & +1.24 \\
Uniform horizon weights & 8/15 & +0.96 \\
Inverse-horizon weights & 8/15 & +14.75 \\
Equal TV weights & 9/15 & +6.26 \\
Adaptive exponent 0.5 & 9/15 & +6.76 \\
Offset $0.001$ & 9/15 & +6.80 \\
Offset $0.1$ & 9/15 & +6.81 \\
Segment penalty $0$ & 9/15 & +6.80 \\
Segment penalty $0.01$ & 9/15 & +6.79 \\
Segment penalty $0.03$ & 9/15 & +6.22 \\
Segment penalty $0.05$ & 9/15 & +6.24 \\
Huber parameter 1 & 9/15 & +6.72 \\
Huber parameter 2 & 9/15 & +6.72 \\
\bottomrule
\end{tabular}
\par\smallskip
\begin{minipage}{0.98\linewidth}\footnotesize
\textit{Notes:} Each row summarizes 15 assets at $H=30$ under expanding estimation relative to the TV experiment's own ridge. The default uses quadratic increment loss, MAD normalization, horizon weights $h^{-1/2}$, adaptive exponent 1, offset 0.01, and segment penalty 0.02. Each alternative changes one setting and recalibrates the TV and persistence weights. Huber parameters refer to training-normalized increment units. These are descriptive results, not MCPB settings or adjusted significance findings.
\end{minipage}
\end{table}
